\chapter{Einleitung}
\label{ch:einleitung}

\section{Motivation und Problemstellung}
\label{sec:motivation}

Dienste hinter NAT oder Firewalls für entfernte Nutzer erreichbar zu machen,
ist ein alltäglicher Bedarf; die verbreiteten Tunneldienste lösen ihn jedoch um
den Preis des Vertrauens: Der Verkehr wird beim Anbieter terminiert, sodass
dieser den Klartext einsehen und die Verbindung auf Zuruf kappen kann. Für
Nutzer mit einem Schutzbedarf gegen Zensur ist ein Anbieter, der den Datenverkehr
\emph{sehen} oder \emph{unterbrechen} kann, jedoch selbst ein Angriffspunkt.

Diese Arbeit verfolgt den entgegengesetzten Entwurf: einen
\emph{providerblinden} (zero-knowledge) Tunnel, bei dem der Betreiber
bewusst aus drei Dingen herausgehalten wird -- dem \emph{Inhalt} (die
Nutzlast-Verschlüsselung terminiert erst am kundeneigenen Origin, nie beim
Betreiber), dem \emph{Vertrauen} (die Schlüssel sind kundenverankert, es gibt
keine betreiberseitige PKI) und, soweit möglich, dem \emph{Datenpfad} selbst.
Die Zero-Knowledge-Eigenschaft ist dabei bewusst auf die \emph{Nutzlast}
begrenzt: Strukturelle Metadaten wie die Existenz eines Tunnels oder -- im
Relay-Fall -- grobe Zeit- und Volumenzähler bleiben sichtbar. Durchsetzung
reduziert sich auf einen einzigen, groben Hebel (Terminierung eines Tokens),
der nur an einer eng definierten \emph{rechtlichen Grenze} greift.

Aus dieser Haltung ergibt sich ein Spannungsfeld: Providerblindheit und
Zensurresistenz erkauft man sich mit zusätzlichen Vermittlungsschritten
(Rendezvous, Relay) und kryptographischem Aufwand. Ob und in welchem Umfang
diese Schutzziele die Latenz eines Tunnels belasten, ist die zentrale empirische
Frage dieser Arbeit.

\section{Zielsetzung}
\label{sec:zielsetzung}

Ziel der Arbeit ist der Entwurf, die prototypische Implementierung und die
emulationsbasierte Evaluation eines solchen Tunnels. Der Prototyp ist in Rust auf
Basis von QUIC umgesetzt und realisiert den providerblinden Relay-Pfad
Client~$\rightarrow$~Edge~$\rightarrow$~Agent~$\rightarrow$~Origin mit einem
Proof-of-Work-gestützten, ratenbegrenzten Rendezvous. Die Noise-basierte
Ende-zu-Ende-Verschlüsselung ist als kryptographischer Baustein umgesetzt. Die
Bewertung erfolgt in einem vollständig Docker-basierten Testbett, in dem
Netzbedingungen (Verzögerung, Paketverlust, Bandbreite) mit \texttt{tc netem}
emuliert und die Roundtrip-Latenz reproduzierbar vermessen werden.

\section{Forschungsfragen}
\label{sec:forschungsfragen}

Die Arbeit gliedert sich an drei Forschungsfragen:

\begin{description}
  \item[FF1 -- Realisierbarkeit.] Lässt sich ein providerblinder,
    zensurresistenter Tunnel aus etablierten Bausteinen -- QUIC als Transport,
    Noise für die Ende-zu-Ende-Sicherheit, Proof-of-Work als Zugangsschutz --
    schlüssig entwerfen und implementieren?
  \item[FF2 -- Latenz-Overhead.] Welchen Grund-Overhead verursacht der
    providerblinde Vermittlungspfad (Verbindungsaufbau, PoW-Rendezvous,
    Mehrsprung-Relay) gegenüber einer direkten Verbindung?
  \item[FF3 -- Verhalten unter Netzbedingungen.] Wie skaliert die
    Tunnel-Roundtrip-Latenz mit Netzverzögerung und Paketverlust?
\end{description}

\section{Aufbau der Arbeit}
\label{sec:aufbau}

Kapitel~\ref{ch:grundlagen} führt die technischen Grundlagen ein
(providerblinde Relays, Noise, QUIC, Proof-of-Work). Kapitel~\ref{ch:architektur}
beschreibt den Systementwurf und dessen Entwurfsentscheidungen.
Kapitel~\ref{ch:implementierung} dokumentiert die Umsetzung des Prototyps.
Kapitel~\ref{ch:evaluation} stellt das Testbett und die Messergebnisse vor und
beantwortet FF2 und FF3. Kapitel~\ref{ch:fazit} fasst die Ergebnisse zusammen
und gibt einen Ausblick auf die offenen Punkte.
